\chapter{El Estado social y su precio}
\label{ch:socialstate}

\section{Lo que se construyó}

Entre 1961 y 1975 el Estado cubano construyó aquello por lo que todavía se le
conoce, y lo hizo lo bastante rápido como para que la literatura internacional
sobre desarrollo lleve discutiéndolo desde entonces. Vale la pena tener la
discusión, pero no debe permitirse que oscurezca los resultados, que no están en
disputa y que se alcanzaron con un ingreso por habitante que nunca superó la
media del rango latinoamericano.

\subsection{Salud}

El punto de partida era malo y empeoró antes de mejorar: alrededor de la mitad de
los médicos del país emigró entre 1959 y 1963, dejando unos tres mil facultativos
para seis millones y medio de personas, la mayoría en La Habana. La respuesta fue
estructural más que paliativa. La formación médica se amplió más allá del punto
de saturación y se mantuvo ahí; se hizo obligatorio el servicio médico rural para
los recién graduados; y el modelo de prestación se reconstruyó en torno al
policlínico, un centro de barrio responsable de una población definida, cuya
tarea principal era la prevención, la vacunación, la atención materna y el
seguimiento de enfermedades crónicas, y no un añadido a la medicina curativa. El
programa del médico de la familia, que puso un médico y una enfermera en cada
cuadra, llegó después, en 1984, pero era la misma lógica bajada un nivel más.

Los resultados medibles fueron grandes. La mortalidad infantil, entre 37 y 40
muertes por cada mil nacidos vivos a finales de los años cincuenta según el
registro incompleto de la época, estaba por debajo de 20 a mediados de los setenta
y llegó a un solo dígito en los ochenta.\unv{} Cuba eliminó la poliomielitis en
1962, primer país de las Américas en hacerlo, mediante una campaña masiva
ejecutada por los comités de cuadra; siguieron el sarampión, la difteria, el
tétanos y la rubéola. La esperanza de vida pasó de unos 62 años a más de 73 a
mediados de los ochenta, convergiendo con la cifra estadounidense y superando a
todos los países de la región salvo Costa Rica \citep{onei,feinsilver1993}.\unv{}

Hay tres matices honestos, y ninguno cancela el logro. Los indicadores de salud
de Cuba ya eran los mejores, o casi los mejores, de América Latina en 1958, de
modo que la revolución partió de una base alta y las ganancias son en parte la
continuación de una tendencia previa ---aunque la convergencia de los resultados
rurales con los urbanos, que es la parte que esa tendencia previa no estaba
entregando, es enteramente posterior a 1959---. La mortalidad infantil es una
estadística sensible a las prácticas de definición y a un sistema de salud con
fuertes incentivos para reducirla; el manejo prenatal agresivo de Cuba, con tasas
muy altas de intervención en los embarazos señalados como de riesgo, contribuye a
la cifra de maneras que una lectura puramente epidemiológica pasa por alto. Y el
logro se compró con una proporción del producto nacional que un país del ingreso
de Cuba solo podía sostener porque otro pagaba el petróleo.

\subsection{Educación}

La escolarización primaria universal se alcanzó de hecho en los años sesenta, y
la secundaria básica universal en los setenta. La matrícula universitaria se
multiplicó aproximadamente por diez en tres décadas. Las escuelas secundarias
internas en el campo combinaron estudio y trabajo agrícola a una escala que
ningún otro país intentó, moviendo a cientos de miles de adolescentes a un
sistema residencial estatal durante parte de cada año. La formación de maestros
se industrializó. Los libros y materiales eran gratuitos.

La calidad era real y se confirmó de forma independiente más tarde. Cuando el
laboratorio regional de la UNESCO evaluó en 1997 a alumnos de tercer y cuarto
grado en trece países latinoamericanos, los niños cubanos obtuvieron alrededor de
350 puntos en una escala cuya media regional era 250 y cuya desviación estándar
era 50 ---aproximadamente dos desviaciones estándar por encima de la región, en
lengua y en matemáticas, con los comparadores más cercanos, Brasil y Chile, cerca
de 260---. El hallazgo que importa más que el promedio es distributivo: los
alumnos de las escuelas cubanas de menor ingreso superaron a la mayoría de los
alumnos de clase media alta del resto de la región \citep{carnoy2007}. Esa es la
cifra más importante para el argumento de la Parte III de este libro, porque es
la que no puede explicarse por el subsidio: dice que el país construyó un sistema
educativo de masas que comprimía las ventajas de origen en vez de reproducirlas,
y que el sistema seguía funcionando después de siete años de Período Especial.

El costo de la educación es que fue también, explícita y deliberadamente, un
sistema de formación política. El currículo llevaba la ideología, el acceso a la
enseñanza superior estaba mediado por el expediente político tanto como por el
académico, y campos enteros ---la sociología, la economía tal como se practica en
otras partes, buena parte de la filosofía--- estuvieron cerrados o reducidos a
catecismo durante largos períodos. Cuba produjo ingenieros, médicos y matemáticos
excelentes, y casi ningún economista, jurista o científico social formado para
pensar cómo funciona un mercado o un tribunal. La Parte III tiene que reparar en
esto: el capital humano del país es magnífico exactamente en las áreas que
necesita una economía técnica, y delgado exactamente en las que necesita una
transición institucional.

\subsection{Las mujeres, y el Código de Familia}

La participación femenina en la fuerza de trabajo pasó de alrededor del 13 por
ciento en 1958 a más del 30 por ciento a fines de los setenta y siguió subiendo,
impulsada por el cuidado infantil gratuito, la educación gratuita y una
Federación de Mujeres Cubanas que era a la vez organización de masas, prestadora
de servicios y correa de transmisión del partido. El Código de Familia de 1975
exigía, en la ley, que los maridos compartieran por igual las tareas domésticas y
el cuidado de los hijos ---disposición inaplicable, y todos sabían que era
inaplicable, y que sin embargo era una declaración que ningún otro país del
hemisferio había hecho---. El aborto fue legal, gratuito y seguro desde 1965,
décadas antes que en la mayor parte de América Latina. El nivel educativo de las
mujeres cubanas superó al de los hombres y sigue por delante.

Lo que no cambió fue la composición del poder. El Buró Político siguió siendo
abrumadoramente masculino durante décadas, y la división doméstica del trabajo
cambió mucho menos de lo que el estatuto decía.

\subsection{La raza, declarada resuelta}

La desegregación de 1959--61 fue rápida, real y aplicada, y es uno de los logros
morales genuinos de la revolución. Lo que siguió fue una decisión de largas
consecuencias: en 1962 la dirigencia declaró abolida la discriminación racial y
cerrada la cuestión racial, y las organizaciones afrocubanas independientes
---incluidas las sociedades y clubes que habían sido la tradición de sociedad
civil negra desde el siglo XIX--- se disolvieron junto con todo lo demás. La raza
desapareció de las categorías censales de análisis y de la discusión pública
permisible durante treinta años.

El efecto fue que una ganancia real quedó consolidada y que todo progreso
ulterior dejó de ser medible. Cuando la raza volvió a la discusión pública cubana
en los años noventa lo hizo porque la economía del dólar había hecho imposible
ignorar la disparidad: las remesas fluían abrumadoramente hacia cubanos blancos,
porque la emigración de los sesenta había sido abrumadoramente blanca, y el
sector turístico contrataba por ``buena presencia''. Cuba entró en su apertura de
mercado con una distribución racial del acceso a la divisa que había pasado
treinta años sin querer medir \citep{delafuente2001}. El
capítulo~\ref{ch:premises} trata esto como una restricción dura y no como una
nota al pie.

\section{Lo que costó}

Los capítulos que describen la política social cubana y los que describen la
represión cubana suelen escribirlos personas distintas. Describen el mismo
Estado, en los mismos años, haciendo ambas cosas con el mismo aparato y a menudo
a través de la misma organización: el comité de cuadra que hacía la campaña de
vacunación era el que llevaba también el expediente de los vecinos.

\subsection{Los campos}

Entre 1965 y 1968 las Unidades Militares de Ayuda a la Producción, las UMAP,
retuvieron a jóvenes en campos de trabajo agrícola en Camagüey bajo disciplina
militar. Las categorías barridas fueron los creyentes religiosos ---testigos de
Jehová, adventistas, católicos practicantes, seminaristas incluidos---, los
homosexuales, y los jóvenes clasificados como antisociales, lo que en la práctica
significaba de pelo largo, sin empleo o aficionados a la música equivocada.
Pasaron por ellos entre veinticinco y treinta y cinco mil personas.\unv{} Hubo
muertes, golpizas y suicidios. Los campos se cerraron en 1968 tras las protestas
de intelectuales cubanos y la presión internacional, y el cierre fue una
admisión, aunque el Estado nunca ha hecho una completa; Castro dijo en 2010 que
la persecución de los homosexuales en aquellos años había sido una ``gran
injusticia'' y que él cargaba con la responsabilidad.

\subsection{El quinquenio gris}

La apertura cultural de comienzos de los sesenta ---que fue real, y que produjo
el instituto de cine, la Casa de las Américas, una cultura seria de revistas
literarias y un arte de cartel que todavía merece mirarse--- terminó en 1971. El
detonante fue el caso del poeta Heberto Padilla, detenido en marzo de 1971 tras
ganar un premio nacional por un libro que la unión de escritores había criticado,
retenido un mes, y liberado después de una autocrítica pública en la que se
denunció a sí mismo y nombró a sus amigos. La escena fue tan evidentemente
coaccionada que quebró la posición de la revolución ante la izquierda europea y
latinoamericana; Sartre, Vargas Llosa, Calvino, Paz, Sontag y decenas más
firmaron en contra, y varios de ellos no volvieron nunca.

El Primer Congreso Nacional de Educación y Cultura, un mes después, oficializó la
línea: el arte era un arma de la revolución, los homosexuales no eran aptos para
trabajar con jóvenes ni para representar a Cuba en el extranjero, y la creencia
religiosa era incompatible con ciertas profesiones. Lo que siguió, de 1971 a 1976
aproximadamente, se conoce en Cuba como el quinquenio gris, y su mecanismo no fue
la detención masiva sino la parametración: escritores, actores, académicos y
maestros eran evaluados contra criterios ideológicos y personales y, si no los
cumplían, apartados de sus profesiones y reasignados a trabajo manual. Se
liquidaron carreras en silencio, por carta administrativa, y muchas nunca se
restituyeron.

\subsection{La cárcel, la salida y el cierre de la puerta}

Los años sesenta vieron una insurgencia rural en el Escambray, combatida entre
1960 y 1965 y conocida oficialmente como la lucha contra bandidos. Fue derrotada,
y la derrota incluyó el traslado forzoso de miles de familias campesinas del
Escambray a nuevos asentamientos en Pinar del Río, en el extremo opuesto de la
isla. La población penal política de los sesenta llegó a las decenas de miles en
su punto máximo; los plantados, los presos que rechazaban el programa de
reeducación y su uniforme, cumplieron condenas de veinte y treinta años.\unv{}
Nunca se ha publicado un conteo.

La emigración continuó y sus condiciones se endurecieron. Cuando Castro abrió el
puerto de Camarioca en 1965, los cubanoamericanos vinieron en barco; la
negociación resultante produjo los Vuelos de la Libertad, dos vuelos chárter
diarios de Varadero a Miami que llevaron a unas 300.000 personas entre 1965 y
1973.\unv{} El régimen le puso precio a usarlos. Solicitar la salida convertía a
una persona en gusano; el solicitante perdía su empleo, era a menudo asignado a
trabajo agrícola durante la espera de años, se le inventariaban y confiscaban los
bienes al partir, y los hombres en edad militar no podían irse en absoluto. Hubo
familias partidas por estas reglas durante décadas. La economía emocional de la
diáspora cubana ---su intensidad, su rechazo al compromiso, su transmisión a lo
largo de tres generaciones--- no se explica del todo sin esta política, y las
propuestas de la Parte III sobre la diáspora tienen que vérselas con un agravio
que es específico y documentado y no difuso.

\subsection{La Ofensiva Revolucionaria de 1968}

En marzo de 1968 el Estado nacionalizó lo que quedaba de pequeño negocio privado:
entre 55.000 y 58.000 establecimientos, esencialmente todos ---bodegas, bares,
puestos de comida, talleres de reparación, lavanderías, vendedores ambulantes,
carpinteros---.\unv{} La razón declarada era que eran parásitos y caldo de
cultivo del mercado negro. El efecto fue abolir la economía cubana de servicios
de la noche a la mañana, entregar sus funciones a un sistema estatal de
distribución sin capacidad alguna para desempeñarlas, y eliminar en una sola ley
a toda la clase de personas que sabía cómo llevar una pequeña empresa.

Es la decisión puramente económica más consecuente del período, y se discute poco
en relación con la reforma agraria porque no involucraba a extranjeros. Cuba pasó
los cincuenta y ocho años siguientes relegalizando periódicamente, en incrementos
estrechos y reversibles, categorías de actividad que eran legales y funcionaban
en 1967: el trabajo por cuenta propia en 1978 y de nuevo en 1993, los
restaurantes privados en 1993, los taxis privados, las habitaciones privadas de
alquiler, los mercados agropecuarios en 1980 y otra vez en 1994, y finalmente, en
2021, las pequeñas y medianas empresas privadas con personalidad jurídica. Cada
uno de esos pasos se presentó como una reforma audaz. Cada uno fue una
restauración parcial de algo que el Estado había abolido en un solo día de marzo
de 1968.

\begin{ledger}{el Estado social, 1961--1975}
\emph{Logrado.} Salud universal y gratuita con la brecha rural-urbana cerrada,
mortalidad infantil más que reducida a la mitad, polio y sarampión eliminados,
esperanza de vida más de diez años arriba. Educación universal y gratuita hasta
el nivel secundario, con resultados medidos dos desviaciones estándar por encima
de la media regional y los mejores resultados en la parte baja de la distribución
social. Participación laboral femenina más que duplicada, cuidado infantil
gratuito, aborto legal, un estatuto de igualdad adelantado para el hemisferio.
Espacio público desegregado por aplicación y no por declaración. Seguridad social
extendida a toda la población.

\emph{Costo.} Campos de trabajo forzado para creyentes, homosexuales e
inconformes, 1965--68. Una purga cultural, 1971--76, ejecutada por reasignación
administrativa y no por detención, y por tanto sin registro ni reparación.
Encarcelamiento político a una escala nunca divulgada. La salida convertida en
delito castigable, con confiscación y años de trabajo forzado como precio de una
solicitud, partiendo familias a lo largo de tres generaciones. Toda empresa
privada por debajo del nivel de la bodega abolida en 1968, destruyendo la
economía de servicios del país y su acervo de destreza empresarial. La raza
declarada resuelta y luego vuelta inmedible durante treinta años.
\end{ledger}
